\section{System Model and Problem Formulation}\label{sec:model}

\subsection{Network and Mobility}
A swarm $\mathcal{U}=\{1,\dots,N\}$ of multirotor UAVs flies at a common altitude $h$ over a
disaster area $\mathcal{A}\subset\mathbb{R}^2$ containing static obstacles. Obstacles taller than
$h$ form a no-fly set $\mathcal{O}\subset\mathcal{A}$ and obstruct radio links; low rubble only
occludes ground targets. A responder radio (node $0$) is fixed at the launch point
$x_0\in\mathcal{A}$; it is the only sink for victim reports and is assumed not to fail. UAV $i$
has position $x_i(t)\in\mathcal{A}\setminus\mathcal{O}$ and speed $\|\dot x_i\|\le v$. Time is
slotted with control period $\Delta t$; every UAV broadcasts one beacon per slot containing its
position, its target ledger and acknowledgements (Section~\ref{sec:method}).

\subsection{Channel Model}
The received SINR of the broadcast from $i$ at $j$ in slot $t$ is
\begin{equation}
\gamma_{ij}(t)=\frac{P_t G^2\, 10^{-(\mathrm{PL}(d_{ij})+n_{ij}L_o+X_{ij})/10}\; |h_{ij}|^2}
{\sigma^2+\sum_{k\in\mathcal{I}_j(t)}\bar P_{kj}},\label{eq:sinr}
\end{equation}
where $\mathrm{PL}(d)=\mathrm{PL}(d_0)+10\eta\log_{10}(d/d_0)$ is log-distance path loss,
$n_{ij}$ the number of obstacles crossing the segment $[x_i,x_j]$ with per-obstruction loss $L_o$,
$X_{ij}\sim\mathcal{N}(0,\sigma_X^2)$ log-normal shadowing evolving with the Gudmundson
correlation $\rho=e^{-\Delta/d_c}$ for endpoint displacement $\Delta$, $|h_{ij}|^2$ unit-mean
Rician fading with factor $K$ ($K=0$, Rayleigh, once obstructed), and $\mathcal{I}_j(t)$ the set
of transmitters colliding in the slot, each present with probability $p_a$. A packet is decoded
iff $\gamma_{ij}\ge\gamma_{\mathrm{th}}$. The link outage probability, averaged over fading and
shadowing, is
\begin{equation}
P_{\mathrm{out}}(d,n)=\mathbb{E}_X\!\left[F_{|h|^2}\!\left(\frac{\gamma_{\mathrm{th}}}{\bar\gamma(d,n)\,10^{X/10}}\right)\right],
\label{eq:pout}
\end{equation}
with $F_{|h|^2}$ the non-central $\chi^2_2$ CDF (computed in closed form; the outer expectation
by Gauss--Hermite quadrature). The \emph{reliable range} for target outage $\varepsilon$ is
$R_\varepsilon(n)=\max\{d: P_{\mathrm{out}}(d,n)\le\varepsilon\}$, and the \emph{reliable graph}
$\mathcal{G}_\varepsilon(t)$ joins $i,j$ iff $P_{\mathrm{out}}(d_{ij},n_{ij})\le\varepsilon$.
Table~\ref{tab:channel} lists the four environment profiles; their parameters follow the ranges
reported for low-altitude aerial and street-level links \cite{tr36777,itup1411}, with
shadowing correlation following \cite{gudmundson1991}; they are not calibrated against measurements
in rubble (Section~\ref{sec:discussion}). For the urban-rubble profile $R_{0.05}(0)=91$\,m and $R_{0.05}(1)=48$\,m,
which places the fixed 50\,m range of the earlier SENTINEL design in context: it is the
reliable range of a link behind one obstruction, not a property of the radio alone.

\begin{table}[t]\centering\caption{Channel profiles ($f=2.4$\,GHz, $P_t=20$\,dBm, $B=20$\,MHz,
$\gamma_{\mathrm{th}}=5$\,dB) and resulting reliable ranges at $\varepsilon=0.05$.}\label{tab:channel}
\begin{tabular}{lcccccc}\toprule
Profile & $\eta$ & $\sigma_X$ & $K$ (LoS) & $L_o$ & $R(0)$ & $R(1)$\\\midrule
open & 2.0 & 2 & 10\,dB & 6\,dB & 2434\,m & 684\,m\\
suburban & 2.5 & 4 & 6\,dB & 10\,dB & 339\,m & 95\,m\\
urban-rubble & 3.0 & 6 & 3\,dB & 6\,dB & 91\,m & 48\,m\\
severe & 3.5 & 8 & Rayleigh & 8\,dB & 35\,m & 21\,m\\\bottomrule
\end{tabular}\end{table}

\subsection{Sensing Model}
Each UAV carries a nadir thermal camera with ground footprint radius $r_f$ and an onboard
detector (Section~\ref{sec:impl}). A look at a cell updates the residual probability that an
undetected victim is present by Bayes' rule, $p\leftarrow p(1-P_d)/(1-pP_d)$ \cite{stone1975}. The detector returns a confidence $c\in[0,1]$ whose distribution depends on
slant range and on view-dependent occlusion (calibrated from the detector, Section~\ref{sec:eval});
false positives from warm debris arrive as a Poisson process per look.

\subsection{Information Model}
Every UAV keeps a ledger of world-frame target entries. Observations within a gate are
associated to one entry; the best confidence per viewing sector is kept and the sectors are fused
into one confidence $c_e$ (Section~\ref{sec:fusion}). An entry is \emph{confirmed} when
$c_e\ge c_{\mathrm{hi}}$ and \emph{tentative} when $c_{\mathrm{lo}}\le c_e<c_{\mathrm{hi}}$.
Ledgers are replicated by gossip with acknowledgements, so each entry carries its holder set
$H_e(t)$. An entry is \emph{secured} once $0\in H_e$ (the responder holds it) and \emph{lost}
if it was confirmed but no alive node holds it.

\subsection{Failure Model}
UAVs fail by crashing (the node and its ledger copy disappear), by radio loss (it flies but
cannot communicate) or by gradual brownout (transmit power decays). We study single,
simultaneous, cascading and Poisson failures with per-UAV rate $\lambda_f$.

\subsection{Problem}
Let $\tau_k$ be the time victim $k$ is first confirmed by any UAV and $T_k\ge\tau_k$ the time
its report reaches the responder ($T_k=\infty$ if never). With mission horizon $T$ and value
discount $\theta$, the team maximises the time-discounted report value
\begin{equation}
J=\frac{1}{|\mathcal{V}|}\sum_{k\in\mathcal{V}} e^{-T_k/\theta},\label{eq:J}
\end{equation}
subject to the dynamics, $x_i\notin\mathcal{O}$, and decentralised information (each UAV acts on
its own sensing and on beacons it actually received). $J$ couples three quantities that prior
work optimises separately: discovery (a victim must be found), delivery (its report must reach
the responder), and timeliness (sooner is better, in the spirit of age of information
\cite{kaul2012aoi}). Maintaining $\mathcal{G}_\varepsilon$ connected at all times makes delivery
immediate but restricts discovery; pure exploration maximises discovery but leaves delivery to
chance contacts, so confirmed information can wait arbitrarily long or be lost with its carrier.
We refer to this tension as the \emph{discovery--delivery trade-off}.
